\usepackage[T1]{fontenc}
\usepackage[utf8]{inputenc}
\usepackage{lmodern}
\usepackage[margin=1.05in]{geometry}
\usepackage{amsmath,amssymb,amsthm,mathtools}
\usepackage{enumitem,booktabs,array}
\usepackage{microtype}
\usepackage{xcolor,graphicx,tikz}
\usepackage{placeins}
\usetikzlibrary{arrows.meta,positioning,calc,fit,backgrounds}
\usepackage{xurl}
\usepackage[colorlinks=true,linkcolor=blue!45!black,
            citecolor=green!35!black,urlcolor=blue!55!black,
            bookmarksnumbered=true]{hyperref}

\ifdefined\pdfglyphtounicode
  \pdfgentounicode=1
  \pdfglyphtounicode{tfm:lmex10/parenleftbig}{0028}
  \pdfglyphtounicode{tfm:lmex10/parenrightbig}{0029}
  \pdfglyphtounicode{tfm:lmex10/braceleftbigg}{007B}
  \pdfglyphtounicode{tfm:lmex10/bracerightbigg}{007D}
  \pdfglyphtounicode{tfm:lmex10/parenleftBigg}{0028}
  \pdfglyphtounicode{tfm:lmex10/parenrightBigg}{0029}
  \pdfglyphtounicode{tfm:lmex10/parenlefttp}{239B}
  \pdfglyphtounicode{tfm:lmex10/parenrighttp}{239E}
  \pdfglyphtounicode{tfm:lmex10/bracelefttp}{23A7}
  \pdfglyphtounicode{tfm:lmex10/braceleftbt}{23A9}
  \pdfglyphtounicode{tfm:lmex10/braceleftmid}{23A8}
  \pdfglyphtounicode{tfm:lmex10/braceex}{23AA}
  \pdfglyphtounicode{tfm:lmex10/parenleftbt}{239D}
  \pdfglyphtounicode{tfm:lmex10/parenrightbt}{23A0}
  \pdfglyphtounicode{tfm:lmex10/parenleftex}{239C}
  \pdfglyphtounicode{tfm:lmex10/parenrightex}{239F}
  \pdfglyphtounicode{tfm:lmex10/angbracketleftBig}{27E8}
  \pdfglyphtounicode{tfm:lmex10/angbracketrightBig}{27E9}
  \pdfglyphtounicode{tfm:lmex10/summationtext}{2211}
  \pdfglyphtounicode{tfm:lmex10/producttext}{220F}
  \pdfglyphtounicode{tfm:lmex10/integraltext}{222B}
  \pdfglyphtounicode{tfm:lmex10/summationdisplay}{2211}
  \pdfglyphtounicode{tfm:lmex10/productdisplay}{220F}
  \pdfglyphtounicode{tfm:lmex10/integraldisplay}{222B}
  \pdfglyphtounicode{tfm:lmex10/hatwide}{02C6}
  \pdfglyphtounicode{tfm:lmex10/radicalbig}{221A}
  \pdfglyphtounicode{tfm:lmsy10/mapsto}{007C}
\fi

\numberwithin{equation}{section}
\newtheorem{theorem}{Theorem}[section]
\newtheorem{proposition}[theorem]{Proposition}
\newtheorem{lemma}[theorem]{Lemma}
\newtheorem{corollary}[theorem]{Corollary}
\theoremstyle{definition}
\newtheorem{definition}[theorem]{Definition}
\newtheorem{notation}[theorem]{Notation}
\newtheorem{convention}[theorem]{Convention}
\theoremstyle{remark}
\newtheorem{remark}[theorem]{Remark}

\newcommand{\R}{\mathbb{R}}
\newcommand{\C}{\mathbb{C}}
\newcommand{\N}{\mathbb{N}}
\newcommand{\Z}{\mathbb{Z}}
\DeclareMathOperator{\Ric}{Ric}
\DeclareMathOperator{\Rm}{Rm}
\DeclareMathOperator{\Scal}{Scal}
\DeclareMathOperator{\Vol}{Vol}
\DeclareMathOperator{\diam}{diam}
\DeclareMathOperator{\dist}{dist}
\DeclareMathOperator{\supp}{supp}
\DeclareMathOperator{\tr}{tr}
\DeclareMathOperator{\Id}{Id}
\DeclareMathOperator{\rank}{rank}

\setlist[enumerate]{leftmargin=*,itemsep=3pt,topsep=5pt}
\setlist[itemize]{leftmargin=*,itemsep=3pt,topsep=5pt}
\setlength{\emergencystretch}{1.5em}
\setcounter{tocdepth}{2}
\hypersetup{pdfauthor={OpenAI},
            pdftitle={Path selection and an elliptic inequality on degenerating Ricci-flat trees},
            pdfsubject={Quasianalytic path selection, Ricci-flat bubble trees, and a renormalized Einstein-Hilbert functional}}
